% Options for packages loaded elsewhere
% Options for packages loaded elsewhere
\PassOptionsToPackage{unicode}{hyperref}
\PassOptionsToPackage{hyphens}{url}
\PassOptionsToPackage{dvipsnames,svgnames,x11names}{xcolor}
%
\documentclass[
  letterpaper,
  DIV=11,
  numbers=noendperiod]{scrartcl}
\usepackage{xcolor}
\usepackage{amsmath,amssymb}
\setcounter{secnumdepth}{-\maxdimen} % remove section numbering
\usepackage{iftex}
\ifPDFTeX
  \usepackage[T1]{fontenc}
  \usepackage[utf8]{inputenc}
  \usepackage{textcomp} % provide euro and other symbols
\else % if luatex or xetex
  \usepackage{unicode-math} % this also loads fontspec
  \defaultfontfeatures{Scale=MatchLowercase}
  \defaultfontfeatures[\rmfamily]{Ligatures=TeX,Scale=1}
\fi
\usepackage{lmodern}
\ifPDFTeX\else
  % xetex/luatex font selection
\fi
% Use upquote if available, for straight quotes in verbatim environments
\IfFileExists{upquote.sty}{\usepackage{upquote}}{}
\IfFileExists{microtype.sty}{% use microtype if available
  \usepackage[]{microtype}
  \UseMicrotypeSet[protrusion]{basicmath} % disable protrusion for tt fonts
}{}
\makeatletter
\@ifundefined{KOMAClassName}{% if non-KOMA class
  \IfFileExists{parskip.sty}{%
    \usepackage{parskip}
  }{% else
    \setlength{\parindent}{0pt}
    \setlength{\parskip}{6pt plus 2pt minus 1pt}}
}{% if KOMA class
  \KOMAoptions{parskip=half}}
\makeatother
% Make \paragraph and \subparagraph free-standing
\makeatletter
\ifx\paragraph\undefined\else
  \let\oldparagraph\paragraph
  \renewcommand{\paragraph}{
    \@ifstar
      \xxxParagraphStar
      \xxxParagraphNoStar
  }
  \newcommand{\xxxParagraphStar}[1]{\oldparagraph*{#1}\mbox{}}
  \newcommand{\xxxParagraphNoStar}[1]{\oldparagraph{#1}\mbox{}}
\fi
\ifx\subparagraph\undefined\else
  \let\oldsubparagraph\subparagraph
  \renewcommand{\subparagraph}{
    \@ifstar
      \xxxSubParagraphStar
      \xxxSubParagraphNoStar
  }
  \newcommand{\xxxSubParagraphStar}[1]{\oldsubparagraph*{#1}\mbox{}}
  \newcommand{\xxxSubParagraphNoStar}[1]{\oldsubparagraph{#1}\mbox{}}
\fi
\makeatother

\usepackage{color}
\usepackage{fancyvrb}
\newcommand{\VerbBar}{|}
\newcommand{\VERB}{\Verb[commandchars=\\\{\}]}
\DefineVerbatimEnvironment{Highlighting}{Verbatim}{commandchars=\\\{\}}
% Add ',fontsize=\small' for more characters per line
\usepackage{framed}
\definecolor{shadecolor}{RGB}{241,243,245}
\newenvironment{Shaded}{\begin{snugshade}}{\end{snugshade}}
\newcommand{\AlertTok}[1]{\textcolor[rgb]{0.68,0.00,0.00}{#1}}
\newcommand{\AnnotationTok}[1]{\textcolor[rgb]{0.37,0.37,0.37}{#1}}
\newcommand{\AttributeTok}[1]{\textcolor[rgb]{0.40,0.45,0.13}{#1}}
\newcommand{\BaseNTok}[1]{\textcolor[rgb]{0.68,0.00,0.00}{#1}}
\newcommand{\BuiltInTok}[1]{\textcolor[rgb]{0.00,0.23,0.31}{#1}}
\newcommand{\CharTok}[1]{\textcolor[rgb]{0.13,0.47,0.30}{#1}}
\newcommand{\CommentTok}[1]{\textcolor[rgb]{0.37,0.37,0.37}{#1}}
\newcommand{\CommentVarTok}[1]{\textcolor[rgb]{0.37,0.37,0.37}{\textit{#1}}}
\newcommand{\ConstantTok}[1]{\textcolor[rgb]{0.56,0.35,0.01}{#1}}
\newcommand{\ControlFlowTok}[1]{\textcolor[rgb]{0.00,0.23,0.31}{\textbf{#1}}}
\newcommand{\DataTypeTok}[1]{\textcolor[rgb]{0.68,0.00,0.00}{#1}}
\newcommand{\DecValTok}[1]{\textcolor[rgb]{0.68,0.00,0.00}{#1}}
\newcommand{\DocumentationTok}[1]{\textcolor[rgb]{0.37,0.37,0.37}{\textit{#1}}}
\newcommand{\ErrorTok}[1]{\textcolor[rgb]{0.68,0.00,0.00}{#1}}
\newcommand{\ExtensionTok}[1]{\textcolor[rgb]{0.00,0.23,0.31}{#1}}
\newcommand{\FloatTok}[1]{\textcolor[rgb]{0.68,0.00,0.00}{#1}}
\newcommand{\FunctionTok}[1]{\textcolor[rgb]{0.28,0.35,0.67}{#1}}
\newcommand{\ImportTok}[1]{\textcolor[rgb]{0.00,0.46,0.62}{#1}}
\newcommand{\InformationTok}[1]{\textcolor[rgb]{0.37,0.37,0.37}{#1}}
\newcommand{\KeywordTok}[1]{\textcolor[rgb]{0.00,0.23,0.31}{\textbf{#1}}}
\newcommand{\NormalTok}[1]{\textcolor[rgb]{0.00,0.23,0.31}{#1}}
\newcommand{\OperatorTok}[1]{\textcolor[rgb]{0.37,0.37,0.37}{#1}}
\newcommand{\OtherTok}[1]{\textcolor[rgb]{0.00,0.23,0.31}{#1}}
\newcommand{\PreprocessorTok}[1]{\textcolor[rgb]{0.68,0.00,0.00}{#1}}
\newcommand{\RegionMarkerTok}[1]{\textcolor[rgb]{0.00,0.23,0.31}{#1}}
\newcommand{\SpecialCharTok}[1]{\textcolor[rgb]{0.37,0.37,0.37}{#1}}
\newcommand{\SpecialStringTok}[1]{\textcolor[rgb]{0.13,0.47,0.30}{#1}}
\newcommand{\StringTok}[1]{\textcolor[rgb]{0.13,0.47,0.30}{#1}}
\newcommand{\VariableTok}[1]{\textcolor[rgb]{0.07,0.07,0.07}{#1}}
\newcommand{\VerbatimStringTok}[1]{\textcolor[rgb]{0.13,0.47,0.30}{#1}}
\newcommand{\WarningTok}[1]{\textcolor[rgb]{0.37,0.37,0.37}{\textit{#1}}}

\usepackage{longtable,booktabs,array}
\usepackage{calc} % for calculating minipage widths
% Correct order of tables after \paragraph or \subparagraph
\usepackage{etoolbox}
\makeatletter
\patchcmd\longtable{\par}{\if@noskipsec\mbox{}\fi\par}{}{}
\makeatother
% Allow footnotes in longtable head/foot
\IfFileExists{footnotehyper.sty}{\usepackage{footnotehyper}}{\usepackage{footnote}}
\makesavenoteenv{longtable}
\usepackage{graphicx}
\makeatletter
\newsavebox\pandoc@box
\newcommand*\pandocbounded[1]{% scales image to fit in text height/width
  \sbox\pandoc@box{#1}%
  \Gscale@div\@tempa{\textheight}{\dimexpr\ht\pandoc@box+\dp\pandoc@box\relax}%
  \Gscale@div\@tempb{\linewidth}{\wd\pandoc@box}%
  \ifdim\@tempb\p@<\@tempa\p@\let\@tempa\@tempb\fi% select the smaller of both
  \ifdim\@tempa\p@<\p@\scalebox{\@tempa}{\usebox\pandoc@box}%
  \else\usebox{\pandoc@box}%
  \fi%
}
% Set default figure placement to htbp
\def\fps@figure{htbp}
\makeatother


% definitions for citeproc citations
\NewDocumentCommand\citeproctext{}{}
\NewDocumentCommand\citeproc{mm}{%
  \begingroup\def\citeproctext{#2}\cite{#1}\endgroup}
\makeatletter
 % allow citations to break across lines
 \let\@cite@ofmt\@firstofone
 % avoid brackets around text for \cite:
 \def\@biblabel#1{}
 \def\@cite#1#2{{#1\if@tempswa , #2\fi}}
\makeatother
\newlength{\cslhangindent}
\setlength{\cslhangindent}{1.5em}
\newlength{\csllabelwidth}
\setlength{\csllabelwidth}{3em}
\newenvironment{CSLReferences}[2] % #1 hanging-indent, #2 entry-spacing
 {\begin{list}{}{%
  \setlength{\itemindent}{0pt}
  \setlength{\leftmargin}{0pt}
  \setlength{\parsep}{0pt}
  % turn on hanging indent if param 1 is 1
  \ifodd #1
   \setlength{\leftmargin}{\cslhangindent}
   \setlength{\itemindent}{-1\cslhangindent}
  \fi
  % set entry spacing
  \setlength{\itemsep}{#2\baselineskip}}}
 {\end{list}}
\usepackage{calc}
\newcommand{\CSLBlock}[1]{\hfill\break\parbox[t]{\linewidth}{\strut\ignorespaces#1\strut}}
\newcommand{\CSLLeftMargin}[1]{\parbox[t]{\csllabelwidth}{\strut#1\strut}}
\newcommand{\CSLRightInline}[1]{\parbox[t]{\linewidth - \csllabelwidth}{\strut#1\strut}}
\newcommand{\CSLIndent}[1]{\hspace{\cslhangindent}#1}



\setlength{\emergencystretch}{3em} % prevent overfull lines

\providecommand{\tightlist}{%
  \setlength{\itemsep}{0pt}\setlength{\parskip}{0pt}}



 


\KOMAoption{captions}{tableheading}
\makeatletter
\@ifpackageloaded{tcolorbox}{}{\usepackage[skins,breakable]{tcolorbox}}
\@ifpackageloaded{fontawesome5}{}{\usepackage{fontawesome5}}
\definecolor{quarto-callout-color}{HTML}{909090}
\definecolor{quarto-callout-note-color}{HTML}{0758E5}
\definecolor{quarto-callout-important-color}{HTML}{CC1914}
\definecolor{quarto-callout-warning-color}{HTML}{EB9113}
\definecolor{quarto-callout-tip-color}{HTML}{00A047}
\definecolor{quarto-callout-caution-color}{HTML}{FC5300}
\definecolor{quarto-callout-color-frame}{HTML}{acacac}
\definecolor{quarto-callout-note-color-frame}{HTML}{4582ec}
\definecolor{quarto-callout-important-color-frame}{HTML}{d9534f}
\definecolor{quarto-callout-warning-color-frame}{HTML}{f0ad4e}
\definecolor{quarto-callout-tip-color-frame}{HTML}{02b875}
\definecolor{quarto-callout-caution-color-frame}{HTML}{fd7e14}
\makeatother
\makeatletter
\@ifpackageloaded{caption}{}{\usepackage{caption}}
\AtBeginDocument{%
\ifdefined\contentsname
  \renewcommand*\contentsname{Table of contents}
\else
  \newcommand\contentsname{Table of contents}
\fi
\ifdefined\listfigurename
  \renewcommand*\listfigurename{List of Figures}
\else
  \newcommand\listfigurename{List of Figures}
\fi
\ifdefined\listtablename
  \renewcommand*\listtablename{List of Tables}
\else
  \newcommand\listtablename{List of Tables}
\fi
\ifdefined\figurename
  \renewcommand*\figurename{Figure}
\else
  \newcommand\figurename{Figure}
\fi
\ifdefined\tablename
  \renewcommand*\tablename{Table}
\else
  \newcommand\tablename{Table}
\fi
}
\@ifpackageloaded{float}{}{\usepackage{float}}
\floatstyle{ruled}
\@ifundefined{c@chapter}{\newfloat{codelisting}{h}{lop}}{\newfloat{codelisting}{h}{lop}[chapter]}
\floatname{codelisting}{Listing}
\newcommand*\listoflistings{\listof{codelisting}{List of Listings}}
\makeatother
\makeatletter
\makeatother
\makeatletter
\@ifpackageloaded{caption}{}{\usepackage{caption}}
\@ifpackageloaded{subcaption}{}{\usepackage{subcaption}}
\makeatother
\usepackage{bookmark}
\IfFileExists{xurl.sty}{\usepackage{xurl}}{} % add URL line breaks if available
\urlstyle{same}
\hypersetup{
  pdftitle={Scikit-longitudinal's TpT splitter: growing decision trees along the time dimension of longitudinal data},
  pdfauthor={Mathias Valla; Simon Provost},
  pdfkeywords={longitudinal data, decision
trees, reproducibility, Python, software},
  colorlinks=true,
  linkcolor={blue},
  filecolor={Maroon},
  citecolor={Blue},
  urlcolor={Blue},
  pdfcreator={LaTeX via pandoc}}


\title{Scikit-longitudinal's TpT splitter: growing decision trees along
the time dimension of longitudinal data}
\usepackage{etoolbox}
\makeatletter
\providecommand{\subtitle}[1]{% add subtitle to \maketitle
  \apptocmd{\@title}{\par {\large #1 \par}}{}{}
}
\makeatother
\subtitle{A short software note for Computo}
\author{Mathias Valla \and Simon Provost}
\date{2026-04-20}
\begin{document}
\maketitle
\begin{abstract}
We present the Time-penalized Trees (TpT) splitter as exposed in
scikit-longitudinal. TpT extends decision-tree split selection to
longitudinal settings by penalizing candidate splits that rely on later
observation waves. This note focuses on software usage and
reproducibility rather than a full methodological treatment. We document
the current TpT API in scikit-longitudinal and provide two compact
real-data examples: a subject-level classification task on the primary
biliary cirrhosis follow-up data and a state-level regression task on
the Produc panel data. The accompanying repository contains a frozen
Python environment, data preparation scripts, and Quarto sources for the
paper.
\end{abstract}


\section{Introduction}\label{introduction}

\texttt{scikit-longitudinal} extends the \texttt{scikit-learn} ecosystem
with estimators and workflows tailored to repeated-measures data
(Provost and Freitas 2025). This software note focuses on the
Time-penalized Trees (TpT) splitter, a decision-tree split rule designed
for longitudinal covariates (Valla 2024). Rather than reproducing the
full methodological development, we document the current TpT API in the
package and provide two compact, fully reproducible real-data examples.
Beyond the original algorithmic presentation, TpT has also been shown to
be a consistent estimator for a non-empty set of observation times
(Valla and Milhaud 2026).

The repository is structured as required by Computo (Computo Team 2020):
the paper is written in Quarto, the environment is frozen in
\texttt{environment.yml}, and the data preparation and analysis code are
kept alongside the manuscript.

\section{\texorpdfstring{The TpT splitter in
\texttt{scikit-longitudinal}}{The TpT splitter in scikit-longitudinal}}\label{the-tpt-splitter-in-scikit-longitudinal}

At a tree node associated with parent time index \(t_p\), TpT scores a
candidate split at time \(t_c\) by multiplying the usual impurity
improvement \(\Delta I\) by an exponential time penalty:

\[
G_{\gamma} = \Delta I \exp \left[-\gamma (t_c - t_p)\right].
\]

The hyperparameter \(\gamma \ge 0\) controls how strongly later
measurements are penalized. When \(\gamma = 0\), TpT reduces to ordinary
impurity-based split selection. As \(\gamma\) increases, the tree
becomes more reluctant to use later waves unless they bring a clearly
stronger impurity reduction.

From a user perspective, the main entry points are
\texttt{TpTDecisionTreeClassifier} and
\texttt{TpTDecisionTreeRegressor}. Both can consume LONG-format data
directly by specifying an identifier column, an observation-time column
and a subject-specific duration column.

For this preliminary software note, we select \(\gamma\) on a small
candidate grid by minimizing the fitted tree's weighted total leaf
impurity on the training split. This is intentionally simple and should
be viewed as a placeholder for a more formal tuning strategy in the
final submission.

{Note to self: a Computo reviewer may reasonably ask for a stronger
justification of this gamma-selection rule, or for a validation-based
alternative. Keep the placeholder wording explicit and avoid overselling
it.}

\begin{tcolorbox}[enhanced jigsaw, titlerule=0mm, bottomrule=.15mm, toptitle=1mm, toprule=.15mm, leftrule=.75mm, title=\textcolor{quarto-callout-note-color}{\faInfo}\hspace{0.5em}{Current scope}, colback=white, colframe=quarto-callout-note-color-frame, colbacktitle=quarto-callout-note-color!10!white, left=2mm, opacityback=0, breakable, bottomtitle=1mm, arc=.35mm, coltitle=black, opacitybacktitle=0.6, rightrule=.15mm]

The present implementation is suitable for compact reproducible examples
and software-oriented demonstrations. Missing-value routing inside the
Cython TpT splitter is still treated as future work, and monotonic
constraints are explicitly disabled for TpT at the wrapper level.

\end{tcolorbox}

\section{Reproducible setup}\label{reproducible-setup}

The repository contains:

\begin{itemize}
\tightlist
\item
  \texttt{scripts/prepare\_example\_data.py} to download and preprocess
  the open datasets;
\item
  \texttt{analysis/tpt\_examples.py} with helper functions used by the
  paper;
\item
  \texttt{environment.yml} with a frozen Python environment.
\end{itemize}

The environment currently pins the TpT-enabled forks by GitHub SHA so
that the examples remain runnable before the upstream releases catch up
with the merged pull requests.

\section{Install, quick start, and main
parameters}\label{install-quick-start-and-main-parameters}

The official \texttt{scikit-longitudinal} getting-started guide is
intentionally lightweight, and the same philosophy fits this note well.
For a standard Python environment, the minimal installation is:

\begin{Shaded}
\begin{Highlighting}[]
\ExtensionTok{pip}\NormalTok{ install scikit{-}longitudinal}
\end{Highlighting}
\end{Shaded}

For LONG-format data, the shortest end-to-end TpT workflow is to pass a
subject identifier, an observation-time column, a subject-level duration
column, and the names of the longitudinal predictors:

\begin{Shaded}
\begin{Highlighting}[]
\ImportTok{from}\NormalTok{ scikit\_longitudinal.estimators.trees }\ImportTok{import}\NormalTok{ TpTDecisionTreeClassifier}

\NormalTok{tpt }\OperatorTok{=}\NormalTok{ TpTDecisionTreeClassifier(}
\NormalTok{    gamma}\OperatorTok{=}\FloatTok{0.01}\NormalTok{,}
\NormalTok{    max\_depth}\OperatorTok{=}\DecValTok{3}\NormalTok{,}
\NormalTok{    min\_samples\_leaf}\OperatorTok{=}\DecValTok{10}\NormalTok{,}
\NormalTok{    assume\_long\_format}\OperatorTok{=}\VariableTok{True}\NormalTok{,}
\NormalTok{    id\_col}\OperatorTok{=}\StringTok{"id"}\NormalTok{,}
\NormalTok{    time\_col}\OperatorTok{=}\StringTok{"time"}\NormalTok{,}
\NormalTok{    duration\_col}\OperatorTok{=}\StringTok{"duration"}\NormalTok{,}
\NormalTok{    time\_step}\OperatorTok{=}\FloatTok{1.0}\NormalTok{,}
\NormalTok{    long\_feature\_columns}\OperatorTok{=}\NormalTok{[}\StringTok{"marker\_a"}\NormalTok{, }\StringTok{"marker\_b"}\NormalTok{],}
\NormalTok{)}

\NormalTok{tpt.fit(X\_long, y\_subject)}
\NormalTok{pred }\OperatorTok{=}\NormalTok{ tpt.predict(X\_long\_test)}
\end{Highlighting}
\end{Shaded}

The small set of parameters that matters most in practice is summarized
in Table~\ref{tbl-tpt-main-parameters}.

\begin{Shaded}
\begin{Highlighting}[]
\NormalTok{pd.DataFrame(}
\NormalTok{    [}
\NormalTok{        \{}
            \StringTok{"parameter"}\NormalTok{: }\StringTok{"gamma"}\NormalTok{,}
            \StringTok{"role"}\NormalTok{: }\StringTok{"Strength of the exponential time penalty. \textasciigrave{}gamma = 0\textasciigrave{} removes the penalty; larger values favour earlier waves."}\NormalTok{,}
\NormalTok{        \},}
\NormalTok{        \{}
            \StringTok{"parameter"}\NormalTok{: }\StringTok{"max\_depth, min\_samples\_leaf"}\NormalTok{,}
            \StringTok{"role"}\NormalTok{: }\StringTok{"Standard tree{-}regularization parameters controlling size and stability of the fitted tree."}\NormalTok{,}
\NormalTok{        \},}
\NormalTok{        \{}
            \StringTok{"parameter"}\NormalTok{: }\StringTok{"assume\_long\_format"}\NormalTok{,}
            \StringTok{"role"}\NormalTok{: }\StringTok{"If \textasciigrave{}True\textasciigrave{}, TpT converts a LONG{-}format dataframe to the wide matrix required by the underlying splitter."}\NormalTok{,}
\NormalTok{        \},}
\NormalTok{        \{}
            \StringTok{"parameter"}\NormalTok{: }\StringTok{"id\_col, time\_col, duration\_col"}\NormalTok{,}
            \StringTok{"role"}\NormalTok{: }\StringTok{"Structural columns needed for the LONG{-}format interface."}\NormalTok{,}
\NormalTok{        \},}
\NormalTok{        \{}
            \StringTok{"parameter"}\NormalTok{: }\StringTok{"time\_step"}\NormalTok{,}
            \StringTok{"role"}\NormalTok{: }\StringTok{"Granularity used to discretize time into wave indices before fitting."}\NormalTok{,}
\NormalTok{        \},}
\NormalTok{        \{}
            \StringTok{"parameter"}\NormalTok{: }\StringTok{"long\_feature\_columns"}\NormalTok{,}
            \StringTok{"role"}\NormalTok{: }\StringTok{"Repeated{-}measures covariates to keep when using the LONG{-}format wrapper API."}\NormalTok{,}
\NormalTok{        \},}
\NormalTok{        \{}
            \StringTok{"parameter"}\NormalTok{: }\StringTok{"features\_group"}\NormalTok{,}
            \StringTok{"role"}\NormalTok{: }\StringTok{"Temporal grouping of wave{-}expanded columns; this is required when fitting on an explicit wide matrix rather than through the LONG{-}format wrapper."}\NormalTok{,}
\NormalTok{        \},}
\NormalTok{    ]}
\NormalTok{)}
\end{Highlighting}
\end{Shaded}

\begin{longtable}[]{@{}lll@{}}

\caption{\label{tbl-tpt-main-parameters}Minimal set of TpT parameters
most relevant for the examples in this note.}

\tabularnewline

\toprule\noalign{}
& parameter & role \\
\midrule\noalign{}
\endhead
\bottomrule\noalign{}
\endlastfoot
0 & gamma & Strength of the exponential time penalty. `gam... \\
1 & max\_depth, min\_samples\_leaf & Standard tree-regularization
parameters contro... \\
2 & assume\_long\_format & If `True`, TpT converts a LONG-format
datafram... \\
3 & id\_col, time\_col, duration\_col & Structural columns needed for
the LONG-format ... \\
4 & time\_step & Granularity used to discretize time into wave ... \\
5 & long\_feature\_columns & Repeated-measures covariates to keep when
usin... \\
6 & features\_group & Temporal grouping of wave-expanded columns;
th... \\

\end{longtable}

When the data are already stored in wave-separated columns, or when TpT
is used inside a broader \texttt{scikit-longitudinal} workflow, the
package documentation recommends the same building blocks as the rest of
the library: a \texttt{LongitudinalPipeline}, a time-aware preprocessing
step such as \texttt{MerWavTimePlus}, and a final estimator. A minimal
pipeline skeleton is:

\begin{Shaded}
\begin{Highlighting}[]
\ImportTok{from}\NormalTok{ scikit\_longitudinal.data\_preparation }\ImportTok{import}\NormalTok{ MerWavTimePlus}
\ImportTok{from}\NormalTok{ scikit\_longitudinal.estimators.trees }\ImportTok{import}\NormalTok{ TpTDecisionTreeClassifier}
\ImportTok{from}\NormalTok{ scikit\_longitudinal.pipeline }\ImportTok{import}\NormalTok{ LongitudinalPipeline}

\NormalTok{pipeline }\OperatorTok{=}\NormalTok{ LongitudinalPipeline(}
\NormalTok{    steps}\OperatorTok{=}\NormalTok{[}
\NormalTok{        (}\StringTok{"merwav"}\NormalTok{, MerWavTimePlus()),}
\NormalTok{        (}\StringTok{"tpt"}\NormalTok{, TpTDecisionTreeClassifier(gamma}\OperatorTok{=}\FloatTok{0.01}\NormalTok{, features\_group}\OperatorTok{=}\NormalTok{feature\_groups)),}
\NormalTok{    ],}
\NormalTok{    features\_group}\OperatorTok{=}\NormalTok{feature\_groups,}
\NormalTok{    non\_longitudinal\_features}\OperatorTok{=}\NormalTok{non\_longitudinal\_features,}
\NormalTok{    feature\_list\_names}\OperatorTok{=}\NormalTok{feature\_names,}
\NormalTok{    update\_feature\_groups\_callback}\OperatorTok{=}\StringTok{"default"}\NormalTok{,}
\NormalTok{)}
\end{Highlighting}
\end{Shaded}

This note keeps the examples at the estimator level, but the API is
compatible with these larger package workflows.

\section{Data sets}\label{data-sets}

The two examples use small public repeated-measures datasets with
complementary temporal structures, summarized in
Table~\ref{tbl-dataset-summary}. We cite \texttt{scikit-longitudinal}
through its JOSS software paper (Provost and Freitas 2025). The
\texttt{pbcseq} data contain follow-up laboratory measurements for the
randomized Mayo Clinic primary biliary cirrhosis trial, with visits
planned at 6 months, 1 year, and yearly thereafter (Therneau, Terry
2025). The \texttt{Produc} data form a balanced annual panel of 48 US
states from 1970 to 1986 (Munnell, Alicia H. 2025).

\begin{Shaded}
\begin{Highlighting}[]
\ImportTok{from}\NormalTok{ analysis.tpt\_examples }\ImportTok{import}\NormalTok{ dataset\_summary\_table}

\NormalTok{dataset\_summary\_table()}
\end{Highlighting}
\end{Shaded}

\begin{verbatim}
/Users/mathiasvalla/Desktop/Projet_TPT/PR SCIKIT-LONG/.venv/lib/python3.10/site-packages/stopit/__init__.py:10: UserWarning: pkg_resources is deprecated as an API. See https://setuptools.pypa.io/en/latest/pkg_resources.html. The pkg_resources package is slated for removal as early as 2025-11-30. Refrain from using this package or pin to Setuptools<81.
  import pkg_resources
\end{verbatim}

\begin{longtable}[]{@{}llllll@{}}

\caption{\label{tbl-dataset-summary}Compact summary of the two public
datasets used in the paper.}

\tabularnewline

\toprule\noalign{}
& dataset & observational\_unit & temporal\_structure & predictors\_used
& target \\
\midrule\noalign{}
\endhead
\bottomrule\noalign{}
\endlastfoot
0 & pbcseq & 312 subjects & 1945 visits total; median 5 visits per
subject... & bili, albumin, protime, ast & binary event by end of
follow-up \\
1 & Produc & 48 states & 17 yearly waves in source data; 16 history
wav... & pc, emp, unemp & final-year gross state product (final\_gsp) \\

\end{longtable}

\section{\texorpdfstring{Example 1: Subject-level classification on
\texttt{pbcseq}}{Example 1: Subject-level classification on pbcseq}}\label{example-1-subject-level-classification-on-pbcseq}

Our classification example uses \texttt{pbcseq}, the sequential primary
biliary cirrhosis follow-up data distributed with the \texttt{survival}
package (Therneau, Terry 2025). We keep four complete longitudinal
biomarkers (\texttt{bili}, \texttt{albumin}, \texttt{protime},
\texttt{ast}), use \texttt{day} as the observation-time variable,
\texttt{futime} as the subject-specific horizon, and define a binary
target equal to one when the subject experiences either transplantation
or death by the end of follow-up. In the processed version used here,
this yields 312 subjects, 1,945 recorded visits, and a median of 5
visits per subject (Table~\ref{tbl-dataset-summary}).

The example is run on a subject-level train/test split, with the
long-format API of \texttt{TpTDecisionTreeClassifier}. The minimal
fit/predict workflow is:

\begin{Shaded}
\begin{Highlighting}[]
\ImportTok{from}\NormalTok{ scikit\_longitudinal.estimators.trees }\ImportTok{import}\NormalTok{ TpTDecisionTreeClassifier}

\NormalTok{feature\_columns }\OperatorTok{=}\NormalTok{ [}\StringTok{"id"}\NormalTok{, }\StringTok{"day"}\NormalTok{, }\StringTok{"futime"}\NormalTok{, }\StringTok{"bili"}\NormalTok{, }\StringTok{"albumin"}\NormalTok{, }\StringTok{"protime"}\NormalTok{, }\StringTok{"ast"}\NormalTok{]}

\NormalTok{clf }\OperatorTok{=}\NormalTok{ TpTDecisionTreeClassifier(}
\NormalTok{    gamma}\OperatorTok{=}\NormalTok{selected\_gamma,}
\NormalTok{    max\_depth}\OperatorTok{=}\DecValTok{3}\NormalTok{,}
\NormalTok{    min\_samples\_leaf}\OperatorTok{=}\DecValTok{10}\NormalTok{,}
\NormalTok{    random\_state}\OperatorTok{=}\DecValTok{42}\NormalTok{,}
\NormalTok{    assume\_long\_format}\OperatorTok{=}\VariableTok{True}\NormalTok{,}
\NormalTok{    id\_col}\OperatorTok{=}\StringTok{"id"}\NormalTok{,}
\NormalTok{    time\_col}\OperatorTok{=}\StringTok{"day"}\NormalTok{,}
\NormalTok{    duration\_col}\OperatorTok{=}\StringTok{"futime"}\NormalTok{,}
\NormalTok{    time\_step}\OperatorTok{=}\FloatTok{365.0}\NormalTok{,}
\NormalTok{    long\_feature\_columns}\OperatorTok{=}\NormalTok{[}\StringTok{"bili"}\NormalTok{, }\StringTok{"albumin"}\NormalTok{, }\StringTok{"protime"}\NormalTok{, }\StringTok{"ast"}\NormalTok{],}
\NormalTok{)}

\NormalTok{clf.fit(train\_df[feature\_columns], train\_df[}\StringTok{"event"}\NormalTok{])}
\NormalTok{pred }\OperatorTok{=}\NormalTok{ clf.predict(test\_df[feature\_columns])}
\NormalTok{proba }\OperatorTok{=}\NormalTok{ clf.predict\_proba(test\_df[feature\_columns])[:, }\DecValTok{1}\NormalTok{]}
\end{Highlighting}
\end{Shaded}

\begin{Shaded}
\begin{Highlighting}[]
\ImportTok{import}\NormalTok{ pandas }\ImportTok{as}\NormalTok{ pd}

\ImportTok{from}\NormalTok{ analysis.tpt\_examples }\ImportTok{import}\NormalTok{ run\_classification\_example}

\NormalTok{pbcseq\_example }\OperatorTok{=}\NormalTok{ run\_classification\_example()}
\NormalTok{pd.DataFrame([pbcseq\_example.metrics]).}\BuiltInTok{round}\NormalTok{(}\DecValTok{3}\NormalTok{)}
\end{Highlighting}
\end{Shaded}

\begin{longtable}[]{@{}lllllllllll@{}}

\caption{\label{tbl-pbcseq-metrics}Subject-level test metrics for the
pbcseq classification example.}

\tabularnewline

\toprule\noalign{}
& train\_subjects & test\_subjects & selected\_gamma & accuracy &
balanced\_accuracy & roc\_auc & tree\_total\_impurity &
tree\_mean\_impurity & tree\_depth & n\_nodes \\
\midrule\noalign{}
\endhead
\bottomrule\noalign{}
\endlastfoot
0 & 234 & 78 & 0.184 & 0.846 & 0.865 & 0.912 & 93.675 & 0.207 & 3 &
13 \\

\end{longtable}

We also report the small hyperparameter set used for this preliminary
example.

\begin{Shaded}
\begin{Highlighting}[]
\NormalTok{pbcseq\_example.hyperparameters}
\end{Highlighting}
\end{Shaded}

\begin{longtable}[]{@{}lll@{}}

\caption{\label{tbl-pbcseq-hyperparameters}Hyperparameter choices for
the pbcseq classification example.}

\tabularnewline

\toprule\noalign{}
& parameter & value \\
\midrule\noalign{}
\endhead
\bottomrule\noalign{}
\endlastfoot
0 & input\_format & long \\
1 & criterion & gini \\
2 & candidate\_gammas & 0.000 to 0.400 by 0.001 \\
3 & selected\_gamma & 0.184 \\
4 & max\_depth & 3 \\
5 & min\_samples\_leaf & 10 \\
6 & time\_step\_days & 365.0 \\
7 & selection\_rule & minimize training tree\_total\_impurity on a
0.0... \\

\end{longtable}

The selected gamma comes from the training-grid search shown in
Table~\ref{tbl-pbcseq-gamma-search} and visualized in
Figure~\ref{fig-pbcseq-gamma-profile}.

\begin{Shaded}
\begin{Highlighting}[]
\NormalTok{pbcseq\_example.gamma\_table.}\BuiltInTok{round}\NormalTok{(}\DecValTok{4}\NormalTok{)}
\end{Highlighting}
\end{Shaded}

\begin{longtable}[]{@{}llllll@{}}

\caption{\label{tbl-pbcseq-gamma-search}Gamma search for the pbcseq
classification example. The search uses a 0.001 grid, but the table
keeps only distinct training tree-total-impurity values and breaks ties
by retaining the smallest gamma.}

\tabularnewline

\toprule\noalign{}
& gamma & tree\_total\_impurity & tree\_mean\_impurity & tree\_depth &
n\_nodes \\
\midrule\noalign{}
\endhead
\bottomrule\noalign{}
\endlastfoot
0 & 0.000 & 204.5029 & 0.4388 & 3 & 9 \\
1 & 0.001 & 174.0299 & 0.4095 & 3 & 11 \\
2 & 0.002 & 221.3713 & 0.4730 & 3 & 11 \\
3 & 0.003 & 166.6620 & 0.3987 & 3 & 9 \\
4 & 0.007 & 161.7476 & 0.3547 & 3 & 11 \\
5 & 0.054 & 194.9955 & 0.3645 & 3 & 11 \\
6 & 0.078 & 153.6755 & 0.2825 & 3 & 11 \\
7 & 0.113 & 94.5752 & 0.2189 & 3 & 11 \\
8 & 0.139 & 111.4667 & 0.2488 & 3 & 11 \\
9 & 0.182 & 228.8063 & 0.3945 & 3 & 13 \\
10 & 0.183 & 147.1358 & 0.2093 & 3 & 13 \\
11 & 0.184 & 93.6752 & 0.2068 & 3 & 13 \\
12 & 0.228 & 135.5551 & 0.1805 & 3 & 13 \\
13 & 0.327 & 147.8284 & 0.1707 & 3 & 13 \\
14 & 0.328 & 119.3562 & 0.1871 & 3 & 11 \\
15 & 0.375 & 130.4875 & 0.1977 & 3 & 11 \\
16 & 0.381 & 132.9564 & 0.2277 & 3 & 13 \\
17 & 0.382 & 140.2477 & 0.2474 & 3 & 13 \\

\end{longtable}

\begin{Shaded}
\begin{Highlighting}[]
\NormalTok{profile }\OperatorTok{=}\NormalTok{ pbcseq\_example.gamma\_table.copy()}
\NormalTok{selected\_gamma }\OperatorTok{=}\NormalTok{ pbcseq\_example.metrics[}\StringTok{"selected\_gamma"}\NormalTok{]}
\NormalTok{selected\_row }\OperatorTok{=}\NormalTok{ profile.loc[profile[}\StringTok{"gamma"}\NormalTok{].eq(selected\_gamma)].iloc[}\DecValTok{0}\NormalTok{]}

\NormalTok{fig, ax }\OperatorTok{=}\NormalTok{ plt.subplots(figsize}\OperatorTok{=}\NormalTok{(}\FloatTok{6.6}\NormalTok{, }\FloatTok{4.1}\NormalTok{))}
\NormalTok{ax.plot(}
\NormalTok{    profile[}\StringTok{"gamma"}\NormalTok{],}
\NormalTok{    profile[}\StringTok{"tree\_total\_impurity"}\NormalTok{],}
\NormalTok{    color}\OperatorTok{=}\NormalTok{PAPER\_COLORS[}\StringTok{"teal"}\NormalTok{],}
\NormalTok{    linewidth}\OperatorTok{=}\FloatTok{2.3}\NormalTok{,}
\NormalTok{    marker}\OperatorTok{=}\StringTok{"o"}\NormalTok{,}
\NormalTok{    markersize}\OperatorTok{=}\FloatTok{4.5}\NormalTok{,}
\NormalTok{    markerfacecolor}\OperatorTok{=}\StringTok{"white"}\NormalTok{,}
\NormalTok{)}
\NormalTok{ax.axvline(selected\_gamma, color}\OperatorTok{=}\NormalTok{PAPER\_COLORS[}\StringTok{"gold"}\NormalTok{], linestyle}\OperatorTok{=}\StringTok{"{-}{-}"}\NormalTok{, linewidth}\OperatorTok{=}\FloatTok{1.5}\NormalTok{)}
\NormalTok{ax.scatter(}
\NormalTok{    [selected\_gamma],}
\NormalTok{    [selected\_row[}\StringTok{"tree\_total\_impurity"}\NormalTok{]],}
\NormalTok{    s}\OperatorTok{=}\DecValTok{70}\NormalTok{,}
\NormalTok{    color}\OperatorTok{=}\NormalTok{PAPER\_COLORS[}\StringTok{"gold"}\NormalTok{],}
\NormalTok{    edgecolor}\OperatorTok{=}\StringTok{"white"}\NormalTok{,}
\NormalTok{    linewidth}\OperatorTok{=}\FloatTok{0.8}\NormalTok{,}
\NormalTok{    zorder}\OperatorTok{=}\DecValTok{3}\NormalTok{,}
\NormalTok{)}
\NormalTok{ax.annotate(}
    \SpecialStringTok{f"selected gamma = }\SpecialCharTok{\{}\NormalTok{selected\_gamma}\SpecialCharTok{:.3f\}}\SpecialStringTok{"}\NormalTok{,}
\NormalTok{    xy}\OperatorTok{=}\NormalTok{(selected\_gamma, selected\_row[}\StringTok{"tree\_total\_impurity"}\NormalTok{]),}
\NormalTok{    xytext}\OperatorTok{=}\NormalTok{(}\DecValTok{14}\NormalTok{, }\DecValTok{12}\NormalTok{),}
\NormalTok{    textcoords}\OperatorTok{=}\StringTok{"offset points"}\NormalTok{,}
\NormalTok{    color}\OperatorTok{=}\NormalTok{PAPER\_COLORS[}\StringTok{"ink"}\NormalTok{],}
\NormalTok{)}
\NormalTok{ax.set\_xlabel(}\StringTok{"Gamma"}\NormalTok{)}
\NormalTok{ax.set\_ylabel(}\StringTok{"Training tree total impurity"}\NormalTok{)}
\NormalTok{ax.margins(x}\OperatorTok{=}\FloatTok{0.04}\NormalTok{)}
\NormalTok{plt.show()}
\end{Highlighting}
\end{Shaded}

\begin{figure}[H]

\centering{

\pandocbounded{\includegraphics[keepaspectratio]{paper_files/figure-latex/fig-pbcseq-gamma-profile-output-1.pdf}}

}

\caption{\label{fig-pbcseq-gamma-profile}Training tree-total-impurity
profile for the pbcseq classification example, after collapsing tied
gamma values.}

\end{figure}%

The resulting subject-level metrics are shown in
Table~\ref{tbl-pbcseq-metrics}. A small preview of the held-out
predictions is given in Table~\ref{tbl-pbcseq-predictions}. With this
finer gamma grid, the same depth-3, \texttt{min\_samples\_leaf=10}
configuration selects \texttt{gamma\ =\ 0.184} and still uses later
history: the learned tree splits at waves 2, 3, 4, and 7 rather than
collapsing to baseline-only rules.

\begin{Shaded}
\begin{Highlighting}[]
\NormalTok{pbcseq\_example.predictions.head(}\DecValTok{8}\NormalTok{).}\BuiltInTok{round}\NormalTok{(}\DecValTok{3}\NormalTok{)}
\end{Highlighting}
\end{Shaded}

\begin{longtable}[]{@{}lllll@{}}

\caption{\label{tbl-pbcseq-predictions}Preview of held-out subject-level
predictions for the pbcseq classification example.}

\tabularnewline

\toprule\noalign{}
& id & observed & predicted & predicted\_probability \\
\midrule\noalign{}
\endhead
\bottomrule\noalign{}
\endlastfoot
22 & 78 & 1 & 1 & 1.000 \\
27 & 105 & 1 & 1 & 1.000 \\
70 & 278 & 1 & 1 & 1.000 \\
73 & 292 & 1 & 1 & 1.000 \\
74 & 293 & 0 & 1 & 1.000 \\
8 & 26 & 1 & 1 & 0.944 \\
26 & 94 & 1 & 1 & 0.944 \\
32 & 117 & 1 & 1 & 0.944 \\

\end{longtable}

The fitted tree is displayed below with the dedicated
\texttt{plot\_tree} exporter, using the wide feature names built
internally by the TpT preprocessing step.

\begin{Shaded}
\begin{Highlighting}[]
\ImportTok{from}\NormalTok{ sklearn }\ImportTok{import}\NormalTok{ tree}
\ImportTok{import}\NormalTok{ matplotlib.pyplot }\ImportTok{as}\NormalTok{ plt}

\NormalTok{fig, ax }\OperatorTok{=}\NormalTok{ plt.subplots(figsize}\OperatorTok{=}\NormalTok{(}\DecValTok{10}\NormalTok{, }\DecValTok{6}\NormalTok{))}
\NormalTok{tree.plot\_tree(}
\NormalTok{    pbcseq\_example.model,}
\NormalTok{    feature\_names}\OperatorTok{=}\NormalTok{pbcseq\_example.feature\_names,}
\NormalTok{    filled}\OperatorTok{=}\VariableTok{True}\NormalTok{,}
\NormalTok{    rounded}\OperatorTok{=}\VariableTok{True}\NormalTok{,}
\NormalTok{    fontsize}\OperatorTok{=}\DecValTok{9}\NormalTok{,}
\NormalTok{    tpt\_time\_scale}\OperatorTok{=}\FloatTok{0.5}\NormalTok{,}
\NormalTok{    ax}\OperatorTok{=}\NormalTok{ax,}
\NormalTok{)}
\NormalTok{ax.set\_title(}\SpecialStringTok{f"TpT classification tree (selected gamma = }\SpecialCharTok{\{}\NormalTok{pbcseq\_example}\SpecialCharTok{.}\NormalTok{metrics[}\StringTok{\textquotesingle{}selected\_gamma\textquotesingle{}}\NormalTok{]}\SpecialCharTok{\}}\SpecialStringTok{)"}\NormalTok{)}
\NormalTok{plt.show()}
\end{Highlighting}
\end{Shaded}

\begin{figure}[H]

\centering{

\pandocbounded{\includegraphics[keepaspectratio]{paper_files/figure-latex/fig-pbcseq-tree-output-1.pdf}}

}

\caption{\label{fig-pbcseq-tree}Selected TpT classification tree on
pbcseq, displayed with the dedicated plot\_tree exporter.}

\end{figure}%

This is precisely the kind of output we want in a software note: the
fitted tree shows explicitly which wave-specific biomarkers were
selected by TpT, while the calling interface remains close to ordinary
\texttt{scikit-learn} usage.

The resulting splits are clinically interpretable: early bilirubin and
prothrombin values appear prominently, which is coherent with their role
as markers of liver dysfunction severity in the PBC setting.

\section{\texorpdfstring{Example 2: State-level regression on
\texttt{Produc}}{Example 2: State-level regression on Produc}}\label{example-2-state-level-regression-on-produc}

Our regression example uses \texttt{Produc}, a public US state
production panel available through Rdatasets (Munnell, Alicia H. 2025).
The source data contain 48 states followed yearly from 1970 to 1986. We
keep three longitudinal predictors available at each wave: private
capital stock (\texttt{pc}), employment (\texttt{emp}), and unemployment
rate (\texttt{unemp}). For each state, we hold out the final recorded
gross state product and predict that \texttt{final\_gsp} from the
preceding 16 yearly observations. This produces a compact, real,
multi-wave regression problem with a richer temporal structure than the
earlier single-feature prototype. To keep the example readable while
still producing a nonzero selected gamma and non-baseline splits on the
training impurity criterion, we use the predictor subset \texttt{pc},
\texttt{emp}, and \texttt{unemp} together with a depth-4 tree,
\texttt{min\_samples\_leaf=1}, and \texttt{min\_samples\_split=2}.

In this example we explicitly use the TpT conversion helper
\texttt{long\_to\_wide(...)}, then fit the regressor on the resulting
wide matrix and its \texttt{feature\_groups} metadata. The minimal
workflow is:

\begin{Shaded}
\begin{Highlighting}[]
\ImportTok{from}\NormalTok{ scikit\_longitudinal.estimators.trees }\ImportTok{import}\NormalTok{ TpTDecisionTreeRegressor}
\ImportTok{from}\NormalTok{ scikit\_longitudinal.estimators.trees.TpT.\_preprocessing }\ImportTok{import}\NormalTok{ long\_to\_wide}

\NormalTok{feature\_columns }\OperatorTok{=}\NormalTok{ [}\StringTok{"state"}\NormalTok{, }\StringTok{"time\_index"}\NormalTok{, }\StringTok{"duration"}\NormalTok{, }\StringTok{"pc"}\NormalTok{, }\StringTok{"emp"}\NormalTok{, }\StringTok{"unemp"}\NormalTok{]}

\NormalTok{wide\_train }\OperatorTok{=}\NormalTok{ long\_to\_wide(}
\NormalTok{    train\_df[feature\_columns],}
\NormalTok{    id\_col}\OperatorTok{=}\StringTok{"state"}\NormalTok{,}
\NormalTok{    time\_col}\OperatorTok{=}\StringTok{"time\_index"}\NormalTok{,}
\NormalTok{    duration\_col}\OperatorTok{=}\StringTok{"duration"}\NormalTok{,}
\NormalTok{    time\_step}\OperatorTok{=}\FloatTok{1.0}\NormalTok{,}
\NormalTok{    feature\_columns}\OperatorTok{=}\NormalTok{[}\StringTok{"pc"}\NormalTok{, }\StringTok{"emp"}\NormalTok{, }\StringTok{"unemp"}\NormalTok{],}
\NormalTok{    target}\OperatorTok{=}\NormalTok{train\_df[}\StringTok{"final\_gsp"}\NormalTok{],}
\NormalTok{)}

\NormalTok{reg }\OperatorTok{=}\NormalTok{ TpTDecisionTreeRegressor(}
\NormalTok{    gamma}\OperatorTok{=}\NormalTok{selected\_gamma,}
\NormalTok{    criterion}\OperatorTok{=}\StringTok{"squared\_error"}\NormalTok{,}
\NormalTok{    max\_depth}\OperatorTok{=}\DecValTok{4}\NormalTok{,}
\NormalTok{    min\_samples\_leaf}\OperatorTok{=}\DecValTok{1}\NormalTok{,}
\NormalTok{    min\_samples\_split}\OperatorTok{=}\DecValTok{2}\NormalTok{,}
\NormalTok{    random\_state}\OperatorTok{=}\DecValTok{42}\NormalTok{,}
\NormalTok{    features\_group}\OperatorTok{=}\NormalTok{wide\_train.feature\_groups,}
\NormalTok{)}

\NormalTok{reg.fit(wide\_train.X, wide\_train.y)}
\NormalTok{pred }\OperatorTok{=}\NormalTok{ reg.predict(wide\_test.X)}
\end{Highlighting}
\end{Shaded}

The next two tables show the same training subjects before and after the
explicit long-to-wide conversion.

\begin{Shaded}
\begin{Highlighting}[]
\ImportTok{from}\NormalTok{ analysis.tpt\_examples }\ImportTok{import}\NormalTok{ run\_regression\_example}

\NormalTok{produc\_example }\OperatorTok{=}\NormalTok{ run\_regression\_example()}
\NormalTok{produc\_example.long\_preview.}\BuiltInTok{round}\NormalTok{(}\DecValTok{2}\NormalTok{)}
\end{Highlighting}
\end{Shaded}

\begin{longtable}[]{@{}lllllllll@{}}

\caption{\label{tbl-produc-long-preview}Preview of the long-format
Produc training data.}

\tabularnewline

\toprule\noalign{}
& state & year & time\_index & duration & pc & emp & unemp &
final\_gsp \\
\midrule\noalign{}
\endhead
\bottomrule\noalign{}
\endlastfoot
0 & ALABAMA & 1970 & 0 & 15 & 35793.80 & 1010.5 & 4.7 & 48409 \\
1 & ALABAMA & 1971 & 1 & 15 & 37299.91 & 1021.9 & 5.2 & 48409 \\
2 & ALABAMA & 1972 & 2 & 15 & 38670.30 & 1072.3 & 4.7 & 48409 \\
3 & ALABAMA & 1973 & 3 & 15 & 40084.01 & 1135.5 & 3.9 & 48409 \\
4 & ALABAMA & 1974 & 4 & 15 & 42057.31 & 1169.8 & 5.5 & 48409 \\
5 & ALABAMA & 1975 & 5 & 15 & 43971.71 & 1155.4 & 7.7 & 48409 \\
6 & ALABAMA & 1976 & 6 & 15 & 50221.57 & 1207.0 & 6.8 & 48409 \\
7 & ALABAMA & 1977 & 7 & 15 & 51084.99 & 1269.2 & 7.4 & 48409 \\

\end{longtable}

\begin{Shaded}
\begin{Highlighting}[]
\NormalTok{produc\_example.wide\_preview.}\BuiltInTok{round}\NormalTok{(}\DecValTok{2}\NormalTok{)}
\end{Highlighting}
\end{Shaded}

\begin{longtable}[]{@{}lllllllll@{}}

\caption{\label{tbl-produc-wide-preview}Preview of the wide-format
matrix produced by long\_to\_wide for the Produc example.}

\tabularnewline

\toprule\noalign{}
& subject\_id & pc\_t0 & pc\_t5 & pc\_t15 & emp\_t0 & emp\_t5 &
unemp\_t5 & final\_gsp \\
\midrule\noalign{}
\endhead
\bottomrule\noalign{}
\endlastfoot
0 & ALABAMA & 35793.80 & 43971.71 & 60688.04 & 1010.5 & 1155.4 & 7.7 &
48409.0 \\
1 & ARIZONA & 23585.99 & 30243.29 & 39246.83 & 547.4 & 729.1 & 12.0 &
46058.0 \\
2 & ARKANSAS & 19749.63 & 23519.27 & 33229.46 & 536.2 & 623.8 & 9.5 &
28168.0 \\
3 & CONNECTICUT & 24082.38 & 30306.10 & 35816.37 & 1197.5 & 1223.4 & 9.1
& 61750.0 \\
4 & DELAWARE & 6108.14 & 7624.03 & 9921.23 & 216.8 & 229.9 & 9.7 &
10072.0 \\
5 & FLORIDA & 57178.05 & 73809.64 & 120719.68 & 2152.1 & 2746.4 & 11.0 &
152269.0 \\
6 & IDAHO & 9202.07 & 11134.59 & 14290.67 & 207.8 & 273.0 & 6.2 &
11672.0 \\
7 & ILLINOIS & 114860.83 & 142558.62 & 164794.50 & 4345.6 & 4418.9 & 7.1
& 183849.0 \\

\end{longtable}

\begin{Shaded}
\begin{Highlighting}[]
\NormalTok{pd.DataFrame([produc\_example.metrics]).}\BuiltInTok{round}\NormalTok{(}\DecValTok{3}\NormalTok{)}
\end{Highlighting}
\end{Shaded}

\begin{longtable}[]{@{}llllllllll@{}}

\caption{\label{tbl-produc-metrics}State-level test metrics for the
Produc regression example.}

\tabularnewline

\toprule\noalign{}
& train\_subjects & test\_subjects & selected\_gamma & mae & r2 &
tree\_total\_impurity & tree\_mean\_impurity & tree\_depth & n\_nodes \\
\midrule\noalign{}
\endhead
\bottomrule\noalign{}
\endlastfoot
0 & 36 & 12 & 0.011 & 25230.961 & 0.86 & 2.603622e+08 & 7232282.78 & 4 &
23 \\

\end{longtable}

\begin{Shaded}
\begin{Highlighting}[]
\NormalTok{produc\_example.hyperparameters}
\end{Highlighting}
\end{Shaded}

\begin{longtable}[]{@{}lll@{}}

\caption{\label{tbl-produc-hyperparameters}Hyperparameter choices for
the Produc regression example.}

\tabularnewline

\toprule\noalign{}
& parameter & value \\
\midrule\noalign{}
\endhead
\bottomrule\noalign{}
\endlastfoot
0 & input\_format & wide (after explicit long\_to\_wide conversion) \\
1 & criterion & squared\_error (MSE) \\
2 & candidate\_gammas & 0.000 to 0.100 by 0.001 \\
3 & selected\_gamma & 0.011 \\
4 & max\_depth & 4 \\
5 & min\_samples\_leaf & 1 \\
6 & min\_samples\_split & 2 \\
7 & time\_step\_years & 1.0 \\
8 & longitudinal\_predictors & pc, emp, unemp \\
9 & target & final\_gsp at the held-out final year \\
10 & selection\_rule & minimize training tree\_total\_impurity on a
0.0... \\

\end{longtable}

\begin{Shaded}
\begin{Highlighting}[]
\NormalTok{produc\_example.gamma\_table.loc[}
\NormalTok{    :,}
\NormalTok{    [}\StringTok{"gamma"}\NormalTok{, }\StringTok{"tree\_total\_impurity"}\NormalTok{, }\StringTok{"root\_time"}\NormalTok{, }\StringTok{"min\_split\_time"}\NormalTok{, }\StringTok{"max\_split\_time"}\NormalTok{],}
\NormalTok{].}\BuiltInTok{round}\NormalTok{(}\DecValTok{3}\NormalTok{)}
\end{Highlighting}
\end{Shaded}

\begin{longtable}[]{@{}llllll@{}}

\caption{\label{tbl-produc-gamma-search}Regression gamma path for the
Produc example on the article grid from 0.000 to 0.100 by 0.001.}

\tabularnewline

\toprule\noalign{}
& gamma & tree\_total\_impurity & root\_time & min\_split\_time &
max\_split\_time \\
\midrule\noalign{}
\endhead
\bottomrule\noalign{}
\endlastfoot
0 & 0.000 & 3.403315e+08 & 8 & 8 & 15 \\
1 & 0.001 & 3.403315e+08 & 1 & 1 & 11 \\
2 & 0.002 & 2.619842e+08 & 1 & 1 & 5 \\
3 & 0.003 & 2.619842e+08 & 1 & 1 & 5 \\
4 & 0.004 & 2.619842e+08 & 1 & 1 & 5 \\
... & ... & ... & ... & ... & ... \\
96 & 0.096 & 3.026745e+08 & 0 & 0 & 0 \\
97 & 0.097 & 3.026745e+08 & 0 & 0 & 0 \\
98 & 0.098 & 3.026745e+08 & 0 & 0 & 0 \\
99 & 0.099 & 3.026745e+08 & 0 & 0 & 0 \\
100 & 0.100 & 3.026745e+08 & 0 & 0 & 0 \\

\end{longtable}

\begin{Shaded}
\begin{Highlighting}[]
\NormalTok{profile }\OperatorTok{=}\NormalTok{ produc\_example.gamma\_table.copy()}
\NormalTok{selected\_gamma }\OperatorTok{=}\NormalTok{ produc\_example.metrics[}\StringTok{"selected\_gamma"}\NormalTok{]}
\NormalTok{selected\_row }\OperatorTok{=}\NormalTok{ profile.loc[profile[}\StringTok{"gamma"}\NormalTok{].eq(selected\_gamma)].iloc[}\DecValTok{0}\NormalTok{]}

\NormalTok{fig, axes }\OperatorTok{=}\NormalTok{ plt.subplots(}\DecValTok{2}\NormalTok{, }\DecValTok{1}\NormalTok{, figsize}\OperatorTok{=}\NormalTok{(}\FloatTok{7.1}\NormalTok{, }\FloatTok{6.4}\NormalTok{), sharex}\OperatorTok{=}\VariableTok{True}\NormalTok{, height\_ratios}\OperatorTok{=}\NormalTok{[}\DecValTok{3}\NormalTok{, }\DecValTok{2}\NormalTok{])}

\NormalTok{ax }\OperatorTok{=}\NormalTok{ axes[}\DecValTok{0}\NormalTok{]}
\NormalTok{ax.plot(}
\NormalTok{    profile[}\StringTok{"gamma"}\NormalTok{],}
\NormalTok{    profile[}\StringTok{"tree\_total\_impurity"}\NormalTok{],}
\NormalTok{    color}\OperatorTok{=}\NormalTok{PAPER\_COLORS[}\StringTok{"coral"}\NormalTok{],}
\NormalTok{    linewidth}\OperatorTok{=}\FloatTok{2.3}\NormalTok{,}
\NormalTok{    marker}\OperatorTok{=}\StringTok{"o"}\NormalTok{,}
\NormalTok{    markersize}\OperatorTok{=}\FloatTok{4.5}\NormalTok{,}
\NormalTok{    markerfacecolor}\OperatorTok{=}\StringTok{"white"}\NormalTok{,}
\NormalTok{)}
\NormalTok{ax.axvline(selected\_gamma, color}\OperatorTok{=}\NormalTok{PAPER\_COLORS[}\StringTok{"gold"}\NormalTok{], linestyle}\OperatorTok{=}\StringTok{"{-}{-}"}\NormalTok{, linewidth}\OperatorTok{=}\FloatTok{1.5}\NormalTok{)}
\NormalTok{ax.scatter(}
\NormalTok{    [selected\_gamma],}
\NormalTok{    [selected\_row[}\StringTok{"tree\_total\_impurity"}\NormalTok{]],}
\NormalTok{    s}\OperatorTok{=}\DecValTok{70}\NormalTok{,}
\NormalTok{    color}\OperatorTok{=}\NormalTok{PAPER\_COLORS[}\StringTok{"gold"}\NormalTok{],}
\NormalTok{    edgecolor}\OperatorTok{=}\StringTok{"white"}\NormalTok{,}
\NormalTok{    linewidth}\OperatorTok{=}\FloatTok{0.8}\NormalTok{,}
\NormalTok{    zorder}\OperatorTok{=}\DecValTok{3}\NormalTok{,}
\NormalTok{)}
\NormalTok{ax.annotate(}
    \SpecialStringTok{f"selected gamma = }\SpecialCharTok{\{}\NormalTok{selected\_gamma}\SpecialCharTok{:.3f\}}\SpecialStringTok{"}\NormalTok{,}
\NormalTok{    xy}\OperatorTok{=}\NormalTok{(selected\_gamma, selected\_row[}\StringTok{"tree\_total\_impurity"}\NormalTok{]),}
\NormalTok{    xytext}\OperatorTok{=}\NormalTok{(}\DecValTok{14}\NormalTok{, }\DecValTok{12}\NormalTok{),}
\NormalTok{    textcoords}\OperatorTok{=}\StringTok{"offset points"}\NormalTok{,}
\NormalTok{    color}\OperatorTok{=}\NormalTok{PAPER\_COLORS[}\StringTok{"ink"}\NormalTok{],}
\NormalTok{)}
\NormalTok{ax.set\_xlabel(}\StringTok{"Gamma"}\NormalTok{)}
\NormalTok{ax.set\_ylabel(}\StringTok{"Training tree total impurity"}\NormalTok{)}
\NormalTok{ax.yaxis.set\_major\_formatter(billion\_formatter)}
\NormalTok{ax.margins(x}\OperatorTok{=}\FloatTok{0.04}\NormalTok{)}

\NormalTok{ax2 }\OperatorTok{=}\NormalTok{ axes[}\DecValTok{1}\NormalTok{]}
\NormalTok{ax2.plot(}
\NormalTok{    profile[}\StringTok{"gamma"}\NormalTok{],}
\NormalTok{    profile[}\StringTok{"root\_time"}\NormalTok{],}
\NormalTok{    color}\OperatorTok{=}\NormalTok{PAPER\_COLORS[}\StringTok{"ink"}\NormalTok{],}
\NormalTok{    linewidth}\OperatorTok{=}\FloatTok{2.0}\NormalTok{,}
\NormalTok{    marker}\OperatorTok{=}\StringTok{"o"}\NormalTok{,}
\NormalTok{    markersize}\OperatorTok{=}\DecValTok{4}\NormalTok{,}
\NormalTok{    label}\OperatorTok{=}\StringTok{"root time"}\NormalTok{,}
\NormalTok{)}
\NormalTok{ax2.plot(}
\NormalTok{    profile[}\StringTok{"gamma"}\NormalTok{],}
\NormalTok{    profile[}\StringTok{"min\_split\_time"}\NormalTok{],}
\NormalTok{    color}\OperatorTok{=}\NormalTok{PAPER\_COLORS[}\StringTok{"teal"}\NormalTok{],}
\NormalTok{    linewidth}\OperatorTok{=}\FloatTok{1.8}\NormalTok{,}
\NormalTok{    marker}\OperatorTok{=}\StringTok{"s"}\NormalTok{,}
\NormalTok{    markersize}\OperatorTok{=}\FloatTok{3.8}\NormalTok{,}
\NormalTok{    label}\OperatorTok{=}\StringTok{"min split time"}\NormalTok{,}
\NormalTok{)}
\NormalTok{ax2.plot(}
\NormalTok{    profile[}\StringTok{"gamma"}\NormalTok{],}
\NormalTok{    profile[}\StringTok{"max\_split\_time"}\NormalTok{],}
\NormalTok{    color}\OperatorTok{=}\NormalTok{PAPER\_COLORS[}\StringTok{"plum"}\NormalTok{],}
\NormalTok{    linewidth}\OperatorTok{=}\FloatTok{1.8}\NormalTok{,}
\NormalTok{    marker}\OperatorTok{=}\StringTok{"\^{}"}\NormalTok{,}
\NormalTok{    markersize}\OperatorTok{=}\DecValTok{4}\NormalTok{,}
\NormalTok{    label}\OperatorTok{=}\StringTok{"max split time"}\NormalTok{,}
\NormalTok{)}
\NormalTok{ax2.axvline(selected\_gamma, color}\OperatorTok{=}\NormalTok{PAPER\_COLORS[}\StringTok{"gold"}\NormalTok{], linestyle}\OperatorTok{=}\StringTok{"{-}{-}"}\NormalTok{, linewidth}\OperatorTok{=}\FloatTok{1.5}\NormalTok{)}
\NormalTok{ax2.set\_xlabel(}\StringTok{"Gamma"}\NormalTok{)}
\NormalTok{ax2.set\_ylabel(}\StringTok{"Selected split time"}\NormalTok{)}
\NormalTok{ax2.legend(frameon}\OperatorTok{=}\VariableTok{False}\NormalTok{, loc}\OperatorTok{=}\StringTok{"upper right"}\NormalTok{, ncol}\OperatorTok{=}\DecValTok{3}\NormalTok{)}
\NormalTok{ax2.margins(x}\OperatorTok{=}\FloatTok{0.04}\NormalTok{)}
\NormalTok{plt.show()}
\end{Highlighting}
\end{Shaded}

\begin{figure}[H]

\centering{

\pandocbounded{\includegraphics[keepaspectratio]{paper_files/figure-latex/fig-produc-gamma-profile-output-1.pdf}}

}

\caption{\label{fig-produc-gamma-profile}Training tree-total-impurity
profile (top) and selected split times (bottom) for the Produc
regression example on the article grid from 0.000 to 0.100 by 0.001.}

\end{figure}%

\begin{Shaded}
\begin{Highlighting}[]
\NormalTok{produc\_example.predictions.head(}\DecValTok{8}\NormalTok{).}\BuiltInTok{round}\NormalTok{(}\DecValTok{2}\NormalTok{)}
\end{Highlighting}
\end{Shaded}

\begin{longtable}[]{@{}lllll@{}}

\caption{\label{tbl-produc-predictions}Preview of held-out state-level
predictions for the Produc regression example.}

\tabularnewline

\toprule\noalign{}
& state & observed & predicted & abs\_error \\
\midrule\noalign{}
\endhead
\bottomrule\noalign{}
\endlastfoot
0 & CALIFORNIA & 464550.0 & 309273.00 & 155277.00 \\
10 & TEXAS & 276313.0 & 309273.00 & 32960.00 \\
2 & GEORGIA & 88827.0 & 67233.12 & 21593.88 \\
11 & VIRGINIA & 88171.0 & 67233.12 & 20937.88 \\
8 & NEW\_JERSEY & 133468.0 & 154405.00 & 20937.00 \\
4 & MICHIGAN & 136433.0 & 154405.00 & 17972.00 \\
9 & SOUTH\_CAROLINA & 39281.0 & 47279.67 & 7998.67 \\
3 & IOWA & 40425.0 & 47279.67 & 6854.67 \\

\end{longtable}

\begin{Shaded}
\begin{Highlighting}[]
\ImportTok{import}\NormalTok{ matplotlib.pyplot }\ImportTok{as}\NormalTok{ plt}
\ImportTok{from}\NormalTok{ sklearn.metrics }\ImportTok{import}\NormalTok{ mean\_absolute\_error, r2\_score}

\NormalTok{palette }\OperatorTok{=}\NormalTok{ \{}
    \StringTok{"Baseline history"}\NormalTok{: PAPER\_COLORS[}\StringTok{"teal"}\NormalTok{],}
    \StringTok{"Mid history"}\NormalTok{: PAPER\_COLORS[}\StringTok{"gold"}\NormalTok{],}
    \StringTok{"Full history"}\NormalTok{: PAPER\_COLORS[}\StringTok{"plum"}\NormalTok{],}
\NormalTok{\}}

\NormalTok{fig, ax }\OperatorTok{=}\NormalTok{ plt.subplots(figsize}\OperatorTok{=}\NormalTok{(}\FloatTok{7.0}\NormalTok{, }\FloatTok{4.8}\NormalTok{))}
\ControlFlowTok{for}\NormalTok{ stage, frame }\KeywordTok{in}\NormalTok{ produc\_example.stage\_predictions.groupby(}\StringTok{"stage"}\NormalTok{, sort}\OperatorTok{=}\VariableTok{False}\NormalTok{):}
\NormalTok{    ax.scatter(}
\NormalTok{        frame[}\StringTok{"observed"}\NormalTok{],}
\NormalTok{        frame[}\StringTok{"predicted"}\NormalTok{],}
\NormalTok{        alpha}\OperatorTok{=}\FloatTok{0.8}\NormalTok{,}
\NormalTok{        label}\OperatorTok{=}\NormalTok{stage,}
\NormalTok{        color}\OperatorTok{=}\NormalTok{palette[stage],}
\NormalTok{        s}\OperatorTok{=}\DecValTok{42}\NormalTok{,}
\NormalTok{        edgecolor}\OperatorTok{=}\StringTok{"white"}\NormalTok{,}
\NormalTok{        linewidth}\OperatorTok{=}\FloatTok{0.5}\NormalTok{,}
\NormalTok{    )}
\NormalTok{limits }\OperatorTok{=}\NormalTok{ [}
\NormalTok{    produc\_example.stage\_predictions[[}\StringTok{"observed"}\NormalTok{, }\StringTok{"predicted"}\NormalTok{]].}\BuiltInTok{min}\NormalTok{().}\BuiltInTok{min}\NormalTok{(),}
\NormalTok{    produc\_example.stage\_predictions[[}\StringTok{"observed"}\NormalTok{, }\StringTok{"predicted"}\NormalTok{]].}\BuiltInTok{max}\NormalTok{().}\BuiltInTok{max}\NormalTok{(),}
\NormalTok{]}
\NormalTok{ax.plot(limits, limits, linestyle}\OperatorTok{=}\StringTok{"{-}{-}"}\NormalTok{, linewidth}\OperatorTok{=}\DecValTok{1}\NormalTok{, color}\OperatorTok{=}\StringTok{"black"}\NormalTok{)}
\NormalTok{ax.set\_xlabel(}\StringTok{"Observed final GSP"}\NormalTok{)}
\NormalTok{ax.set\_ylabel(}\StringTok{"Predicted final GSP"}\NormalTok{)}
\NormalTok{ax.set\_title(}\StringTok{"TpT regression on Produc"}\NormalTok{)}
\NormalTok{ax.xaxis.set\_major\_formatter(thousands\_formatter)}
\NormalTok{ax.yaxis.set\_major\_formatter(thousands\_formatter)}
\NormalTok{ax.legend(frameon}\OperatorTok{=}\VariableTok{False}\NormalTok{, loc}\OperatorTok{=}\StringTok{"upper left"}\NormalTok{)}

\NormalTok{stage\_text }\OperatorTok{=} \StringTok{"}\CharTok{\textbackslash{}n}\StringTok{"}\NormalTok{.join(}
\NormalTok{    [}
        \SpecialStringTok{f"}\SpecialCharTok{\{}\NormalTok{stage}\SpecialCharTok{\}}\SpecialStringTok{: MAE=}\SpecialCharTok{\{}\NormalTok{mean\_absolute\_error(frame[}\StringTok{\textquotesingle{}observed\textquotesingle{}}\NormalTok{], frame[}\StringTok{\textquotesingle{}predicted\textquotesingle{}}\NormalTok{])}\SpecialCharTok{:.0f\}}\SpecialStringTok{, R\^{}2=}\SpecialCharTok{\{}\NormalTok{r2\_score(frame[}\StringTok{\textquotesingle{}observed\textquotesingle{}}\NormalTok{], frame[}\StringTok{\textquotesingle{}predicted\textquotesingle{}}\NormalTok{])}\SpecialCharTok{:.2f\}}\SpecialStringTok{"}
        \ControlFlowTok{for}\NormalTok{ stage, frame }\KeywordTok{in}\NormalTok{ produc\_example.stage\_predictions.groupby(}\StringTok{"stage"}\NormalTok{, sort}\OperatorTok{=}\VariableTok{False}\NormalTok{)}
\NormalTok{    ]}
\NormalTok{)}
\NormalTok{ax.text(}
    \FloatTok{0.98}\NormalTok{,}
    \FloatTok{0.03}\NormalTok{,}
\NormalTok{    stage\_text,}
\NormalTok{    transform}\OperatorTok{=}\NormalTok{ax.transAxes,}
\NormalTok{    ha}\OperatorTok{=}\StringTok{"right"}\NormalTok{,}
\NormalTok{    va}\OperatorTok{=}\StringTok{"bottom"}\NormalTok{,}
\NormalTok{    fontsize}\OperatorTok{=}\FloatTok{8.5}\NormalTok{,}
\NormalTok{    color}\OperatorTok{=}\NormalTok{PAPER\_COLORS[}\StringTok{"ink"}\NormalTok{],}
\NormalTok{    bbox}\OperatorTok{=}\NormalTok{\{}\StringTok{"boxstyle"}\NormalTok{: }\StringTok{"round,pad=0.35"}\NormalTok{, }\StringTok{"facecolor"}\NormalTok{: }\StringTok{"white"}\NormalTok{, }\StringTok{"edgecolor"}\NormalTok{: PAPER\_COLORS[}\StringTok{"grid"}\NormalTok{]\},}
\NormalTok{)}
\NormalTok{plt.show()}
\end{Highlighting}
\end{Shaded}

\begin{figure}[H]

\centering{

\pandocbounded{\includegraphics[keepaspectratio]{paper_files/figure-latex/fig-produc-predictions-output-1.pdf}}

}

\caption{\label{fig-produc-predictions}Observed versus predicted final
gross state product in the Produc regression example, with predictions
computed from baseline-only, mid-history, and full-history inputs.}

\end{figure}%

The selected regression tree can also be shown directly with the
dedicated \texttt{plot\_tree} exporter.

On the article grid from \texttt{0.000} to \texttt{0.100} by
\texttt{0.001}, the selected regression gamma is \texttt{0.011}. The
corresponding tree still uses later history, with selected split times
between waves 1 and 5. As gamma increases further on that grid, the tree
is pulled back toward baseline: the root split moves to wave 0 by about
\texttt{gamma\ =\ 0.031}, and from about \texttt{gamma\ =\ 0.055} onward
all selected split times are at baseline.

\begin{Shaded}
\begin{Highlighting}[]
\NormalTok{fig, ax }\OperatorTok{=}\NormalTok{ plt.subplots(figsize}\OperatorTok{=}\NormalTok{(}\DecValTok{13}\NormalTok{, }\FloatTok{6.6}\NormalTok{))}
\NormalTok{tree.plot\_tree(}
\NormalTok{    produc\_example.model,}
\NormalTok{    feature\_names}\OperatorTok{=}\NormalTok{produc\_example.feature\_names,}
\NormalTok{    filled}\OperatorTok{=}\VariableTok{True}\NormalTok{,}
\NormalTok{    rounded}\OperatorTok{=}\VariableTok{True}\NormalTok{,}
\NormalTok{    fontsize}\OperatorTok{=}\DecValTok{7}\NormalTok{,}
\NormalTok{    tpt\_time\_scale}\OperatorTok{=}\FloatTok{0.18}\NormalTok{,}
\NormalTok{    ax}\OperatorTok{=}\NormalTok{ax,}
\NormalTok{)}
\NormalTok{ax.set\_title(}\SpecialStringTok{f"TpT regression tree (selected gamma = }\SpecialCharTok{\{}\NormalTok{produc\_example}\SpecialCharTok{.}\NormalTok{metrics[}\StringTok{\textquotesingle{}selected\_gamma\textquotesingle{}}\NormalTok{]}\SpecialCharTok{\}}\SpecialStringTok{)"}\NormalTok{)}
\NormalTok{plt.show()}
\end{Highlighting}
\end{Shaded}

\begin{figure}[H]

\centering{

\pandocbounded{\includegraphics[keepaspectratio]{paper_files/figure-latex/fig-produc-tree-output-1.pdf}}

}

\caption{\label{fig-produc-tree}Selected TpT regression tree on Produc,
displayed with the dedicated plot\_tree exporter.}

\end{figure}%

{Note to self: Produc is a much better multi-wave regression example
than the earlier single-feature draft, but it is still panel-economic
data rather than an individual-level biomedical cohort. If a stronger
public longitudinal regression dataset appears, keep replacement open.}

{Note to self: the visual style is still functional rather than
polished. Before submission, improve the final typography and harmonize
colors across tables, tree plots, and scatter plots.}

\section{Discussion}\label{discussion}

These two examples are intentionally small. Their role is not to provide
a full benchmarking study, but to document a reproducible and
inspectable usage path for TpT within \texttt{scikit-longitudinal}.
Together they illustrate four software-facing points:

\begin{enumerate}
\def\labelenumi{\arabic{enumi}.}
\tightlist
\item
  subject-level train/test evaluation can be carried out while keeping
  the inputs in LONG format, or by explicit conversion to the wide
  matrix expected by the core TpT splitter;
\item
  TpT exposes a familiar \texttt{scikit-learn}-like interface through
  \texttt{fit}, \texttt{predict}, and \texttt{predict\_proba};
\item
  preliminary gamma tuning can be implemented with a compact
  impurity-based helper before moving to a more formal model-selection
  strategy;
\item
  the learned trees remain readable because the selected splits are tied
  to explicit wave-indexed feature names.
\end{enumerate}

Most importantly for this note, the exported trees show both what
matters and when it matters. In \texttt{pbcseq}, the node labels tell us
whether bilirubin or prothrombin is selected from an early or later
visit; in \texttt{Produc}, they reveal whether public capital, private
capital, employment, or unemployment becomes decisive at the beginning
or later in the state-level history. This explicit predictor-plus-wave
readability is one of the main practical strengths of TpT.

Future revisions of this note can stay short while becoming more robust
by switching the environment from pinned GitHub SHAs to released package
versions, and by updating the examples once the remaining TpT
implementation limitations around missing-value routing are resolved.

{Note to self: explicitly state near submission time whether the paper
is tied to released package versions or to pinned development SHAs,
because Computo reviewers will care about long-term reproducibility.}

\section{Acknowledgements}\label{acknowledgements}

Work(s) conducted within the Research Chair DIALog under the aegis of
the Risk Foundation, a joint initiative by Université Paris-Dauphine PSL
and CNP Assurances.

Placeholder for Simon Provost's acknowledgements.

{Note to self: fill Simon's affiliation and acknowledgements before
submission, and check whether Computo wants a software-availability or
funding sentence in a more specific location.}

\section*{References}\label{references}
\addcontentsline{toc}{section}{References}

{Note to self: complete and clean the bibliography before submission,
including any missing software, package, and methodological references
cited in the final text.}

\phantomsection\label{refs}
\begin{CSLReferences}{1}{0}
\bibitem[\citeproctext]{ref-computo}
Computo Team. 2020. {``Computo: Reproducible Computational/Algorithmic
Contributions in Statistics and Machine Learning.''}

\bibitem[\citeproctext]{ref-dataset-produc}
Munnell, Alicia H. 2025. {``Produc: US State Production Panel.''}
\url{https://vincentarelbundock.github.io/Rdatasets/doc/Ecdat/Produc.html}.

\bibitem[\citeproctext]{ref-provost2025sklong}
Provost, Simon, and Alex A. Freitas. 2025. {``Scikit-Longitudinal: A
Machine Learning Library for Longitudinal Classification in Python.''}
\emph{Journal of Open Source Software} 10 (112): 8481.
\url{https://doi.org/10.21105/joss.08481}.

\bibitem[\citeproctext]{ref-dataset-pbcseq}
Therneau, Terry. 2025. {``Pbcseq: Mayo Clinic Primary Biliary Cirrhosis
Sequential Data.''} \url{https://rdrr.io/cran/survival/man/pbcseq.html}.

\bibitem[\citeproctext]{ref-valla2024tpt}
Valla, Mathias. 2024. {``Time-Penalized Trees ({TpT}): Introducing a New
Tree-Based Data Mining Algorithm for Time-Varying Covariates.''}
\emph{Annals of Mathematics and Artificial Intelligence} 92: 1609--61.
\url{https://doi.org/10.1007/s10472-024-09950-w}.

\bibitem[\citeproctext]{ref-valla2026consistency}
Valla, Mathias, and Xavier Milhaud. 2026. {``Consistent Time-Aware Trees
for Longitudinal Data: The Time-Penalized Tree.''}
\url{https://cnrs.hal.science/hal-05022929}.

\end{CSLReferences}




\end{document}
